\documentclass[12pt,a4paper]{amsart}

\usepackage{amsmath,amssymb,amsthm}
\usepackage{mathtools}
\usepackage{enumerate}
\usepackage{hyperref}
\usepackage{cleveref}
\usepackage{geometry}
\geometry{margin=2.5cm}

%%------------------------------------------------------------------
%% Theorem environments
%%------------------------------------------------------------------
\newtheorem{theorem}{Theorem}[section]
\newtheorem{lemma}[theorem]{Lemma}
\newtheorem{corollary}[theorem]{Corollary}
\newtheorem{proposition}[theorem]{Proposition}
\newtheorem{definition}[theorem]{Definition}
\newtheorem{remark}[theorem]{Remark}
\newtheorem{assumption}[theorem]{Assumption}

%%------------------------------------------------------------------
%% Notation macros
%%------------------------------------------------------------------
\newcommand{\psin}{\psi_n}
\newcommand{\wncl}{\omega_n^{\mathrm{cl}}}
\newcommand{\rhocl}{\rho^{\mathrm{cl}}}
\newcommand{\IB}{I_\delta}

\begin{document}

\title[Semiclassical Equidistribution of PSWF Densities]%
{Semiclassical Equidistribution of Prolate Spheroidal Wave Function
Densities via Pr\"ufer Phase Analysis}

\author{Anonymous}
\date{April 2026}
\maketitle

%%=================================================================
\begin{abstract}
We prove that the squared densities $\psin^2$ of the prolate
spheroidal wave functions (PSWFs) satisfy a uniform semiclassical
approximation against the classical density $\rhocl_n$ on $[-T,T]$
at a rate $O(1/n)$ for all $n \le \gamma N$, where $N = cT/\pi$
is the Shannon number and the asymptotic regime is $N \to \infty$
(equivalently $c \to \infty$ with $T$ fixed).
The proof uses Pr\"ufer phase analysis, integration by parts for
the oscillatory terms, and an explicit drift cancellation identity
that reduces the amplitude deviation to a single oscillatory
integral. No microlocal machinery, Airy-function matching, or
Riemann--Hilbert methods are used.
This result supplies the semiclassical equidistribution hypothesis
required by the bulk tail bound program of \cite{PaperIII}.
\end{abstract}

%%=================================================================
\section{Introduction}\label{sec:intro}
%%=================================================================

\subsection{Background and motivation}

Prolate spheroidal wave functions $\psin$ are the eigenfunctions
of the finite Fourier operator $\mathcal{F}_T$ on $[-T,T]$:
\[
  \int_{-T}^T \psin(y)\,e^{icxy}\,dy = \mu_n(c)\,\psin(x),
\]
where $c > 0$ is the bandwidth parameter and
$\lambda_n(c) = |\mu_n(c)|^2$ are the eigenvalues.
Equivalently, $\psin$ are eigenfunctions of the differential
operator
\[
  \mathcal{L}_c\psin
  = \bigl[(1-x^2/T^2)\psin'\bigr]'
    + \bigl[c^2(1-x^2/T^2) - \chi_n(c)/T^2\bigr]\psin = 0,
\]
with eigenvalues $\chi_n(c)$ satisfying
$\chi_n(c) \sim n(n+1) + (2n+1)c/\pi + O(c^2/n)$
for $n$ in the bulk regime $n \le \gamma N$, $N = cT/\pi$.

The present paper establishes a quantitative semiclassical
equidistribution theorem for $\psin^2$, which is a key analytic
ingredient in the bulk tail bound program developed in
\cite{PaperI,PaperII,PaperIII}.

\subsection{Main result (informal)}

For $n$ in the bulk regime $n \le \gamma N$, the squared density
$\psin^2$, when renormalized by $\lambda_n(c)$, satisfies the
uniform semiclassical approximation against the classical
reference density
\[
  \rhocl_n(x) =
  \frac{1/\wncl(x)}{\int_{-x_+(n)}^{x_+(n)}
                  dt/\wncl(t)},
  \quad |x| < x_+(n),
\]
at rate $O(1/n)$, uniformly over $C^1$ test functions and
$n \le \gamma N$.

\begin{remark}[Nature of the result]\label{rem:not-weak-conv}
Theorem~\ref{thm:weak-limit} is \emph{not} a convergence theorem
in the sense of weak convergence to a fixed limiting measure.
Instead, it is a \emph{uniform semiclassical matching principle}
between the spectral density $\psin^2$ and the $n$-dependent
classical reference density $\rhocl_n$.
In particular, $\rhocl_n$ depends on $n$ through the turning
point $x_+(n)$ and should be understood as the
\emph{local semiclassical matching density} for index $n$,
not a fixed target distribution.
\end{remark}

\subsection{Strategy of proof}

The proof has four layers, executed in Sections
\ref{sec:prufer}--\ref{sec:weaklimit}:

\begin{enumerate}[(i)]
  \item \textbf{Pr\"ufer transform} (Section~\ref{sec:prufer}):
  reduce $\psin^2$ to amplitude $r_n^2$ and phase $\theta_n$.
  \item \textbf{Oscillation control}
  (Section~\ref{sec:oscillation}):
  show $\int h\cos(2\theta_n) = O(\|h\|_{C^1}/n)$ by IBP,
  using the monotonicity $\theta_n' \ge c_{\delta,\gamma}n > 0$
  on $I_\delta$.
  \item \textbf{Amplitude control}
  (Section~\ref{sec:amplitude}):
  establish pointwise comparability
  $r_n(x)^2 \asymp r_n(x_0)^2$
  (Lemma~\ref{lem:pointwise-control}),
  then show $r_n^2/r_n^{\mathrm{WKB},2} = 1 + O(1/n)$
  via the drift cancellation identity and a direct
  Gronwall estimate (Lemma~\ref{lem:amplitude-drift}).
  \item \textbf{Main theorem} (Section~\ref{sec:weaklimit}):
  assemble the above to obtain the rate-$O(1/n)$
  equidistribution.
\end{enumerate}

\subsection{Logical structure and non-circularity}

\begin{center}
\begin{tabular}{lll}
\hline
Step & Uses & Does \emph{not} use \\
\hline
Lemma~\ref{lem:prufer-oscillation-control} &
  $\theta_n' \ge c_{\delta,\gamma}n$ &
  $r_n$, $\Delta_n$ \\
Lemma~\ref{lem:deviation-ode} &
  Pr\"ufer ODE &
  Lemma~\ref{lem:prufer-oscillation-control} \\
Lemma~\ref{lem:pointwise-control} &
  Pr\"ufer ODE,
  Lemmas~\ref{lem:G-regularity}, \ref{lem:phase-lower},
  \ref{lem:phase-upper} &
  Lemmas~\ref{lem:amplitude-drift},
  \ref{lem:normalization},
  Theorem~\ref{thm:weak-limit} \\
Lemma~\ref{lem:amplitude-drift} &
  Lemmas~\ref{lem:prufer-oscillation-control},
  \ref{lem:deviation-ode},
  \ref{lem:pointwise-control} &
  Theorem~\ref{thm:weak-limit} \\
Lemma~\ref{lem:normalization} &
  Lemma~\ref{lem:amplitude-drift} &
  Theorem~\ref{thm:weak-limit} \\
Theorem~\ref{thm:weak-limit} &
  Lemmas~\ref{lem:prufer-oscillation-control},
  \ref{lem:amplitude-drift},
  \ref{lem:normalization} &
  --- \\
\hline
\end{tabular}
\end{center}

\noindent
In particular, no step uses Theorem~\ref{thm:weak-limit} as a
hypothesis; the dependency graph is a strict DAG with no cycles.

\subsection{Relation to prior work}

Oscillatory estimates for PSWF densities appear implicitly in
\cite{OsipovRokhlinXiao}, but without quantitative rates.
The Pr\"ufer-based approach is classical for Sturm--Liouville
theory \cite{Levitan}; its adaptation to the PSWF setting
with explicit $O(1/n)$ rates is new.
The result is a PSWF analogue of the Szeg\H{o} equidistribution
theorem for orthogonal polynomial measures \cite{Szego},
proved here by elementary ODE methods rather than
potential-theoretic or Riemann--Hilbert techniques.

%%=================================================================
\section{Setup, Notation, and Preliminary Lemmas}\label{sec:setup}
%%=================================================================

\subsection{Parameters and bulk regime}

Throughout, $c, T > 0$ are fixed reference values,
$N := cT/\pi$ is the Shannon number, and
$0 < \gamma < 1$ is a fixed bulk parameter.
We write $n \le \gamma N$ for the bulk regime.

\subsection{Asymptotic regime}\label{subsec:regime}

Throughout the paper we work in the semiclassical bulk regime
\begin{equation}\label{eq:semiclassical-regime}
  N \to \infty,
  \qquad \gamma \in (0,1)\ \text{fixed},
  \qquad n \le \gamma N,
\end{equation}
where $N = cT/\pi$.
Equivalently, this corresponds to the semiclassical limit
\[
  c \to \infty \quad \text{with } T > 0\ \text{fixed},
\]
so that the bandwidth parameter $c$ acts as the large parameter
of the problem.

All estimates of the form $O(1/n)$ are therefore understood
in the joint limit $n, N \to \infty$ with $n/N \le \gamma < 1$.
In particular, all constants are uniform in this regime and
depend only on $\delta$, $\gamma$, $c$, and $T$
(with $c$ viewed as the scaling parameter generating the limit).

\begin{remark}[Meaning of uniformity]\label{rem:uniformity}
All bounds stated as $O(\cdot)$ or $\le C_{\delta,\gamma}(\cdots)$
hold \emph{uniformly for all $n \le \gamma N$} in the regime
\eqref{eq:semiclassical-regime}.
The constants $C_{\delta,\gamma}$, $c_{\delta,\gamma}$
depend only on $\delta$, $\gamma$, $c$, $T$, and are independent
of the individual value of $n$.
We write $C_{\delta,\gamma}$ throughout as shorthand
for constants that may also depend on the fixed parameters
$c$ and $T$.
\end{remark}

\subsection{Turning points and classical frequency}

For $n \le \gamma N$, the classical turning point is:
\begin{equation}\label{eq:turning-point}
  x_+(n) := \frac{T}{c}\sqrt{\chi_n(c)},
\end{equation}
with $x_+(n) \le (1-\eta_\gamma)T$ for some $\eta_\gamma > 0$
depending only on $\gamma$.

The classical frequency function is:
\begin{equation}\label{eq:wncl-def}
  \wncl(x) := \sqrt{c^2(1 - x^2/T^2) + \chi_n(c)/T^2},
\end{equation}
which satisfies $\wncl(x) \asymp n$ uniformly on $I_\delta$
(see Lemma~\ref{lem:G-regularity} below).

\subsection{Bulk interval}

Fix $0 < \delta < \eta_\gamma$. The parameter $\delta > 0$ is the
\emph{bulk separation parameter}: points $x \in I_\delta$ lie at
distance at least $\delta T$ from the endpoints $\pm T$, with all
constants depending on $\delta$ but uniform in $n \le \gamma N$.
Define the bulk interval:
\begin{equation}\label{eq:Idelta}
  I_\delta := [-(1-\delta)T,\,(1-\delta)T].
\end{equation}
Then $[-x_+(n), x_+(n)] \subset I_\delta \subset [-T,T]$
for all $n \le \gamma N$.

\subsection{Classical density}

\begin{definition}[Classical density]\label{def:rho-cl}
\begin{equation}\label{eq:rho-cl}
  \rhocl_n(x) :=
  \frac{1/\wncl(x)}{\int_{-x_+(n)}^{x_+(n)} dt/\wncl(t)},
  \quad x \in [-x_+(n), x_+(n)].
\end{equation}
The normalization $\int_{-x_+(n)}^{x_+(n)}\rhocl_n(x)\,dx = 1$
holds by definition.
Note that $\rhocl_n$ depends on $n$ through the turning point
$x_+(n)$; it is the \emph{local semiclassical matching density}
for index $n$, not a fixed target distribution.
By Lemma~\ref{lem:denom-bs-pswf} below,
the denominator satisfies
\begin{equation}\label{eq:denom-BS}
  \int_{-x_+(n)}^{x_+(n)}\frac{dt}{\wncl(t)}
  = \frac{\pi T}{n + \tfrac{1}{2}} + O\!\left(\frac{1}{n^2}\right),
\end{equation}
where the $O(1/n^2)$ relative error propagates
to an $O(1/n)$ relative error in $\rhocl_n$ only through the
normalization constant; no mass is lost in the cutoff
$I_\delta \to [-x_+(n),x_+(n)]$ since
$\mathrm{supp}\,\rhocl_n = [-x_+(n),x_+(n)] \subset I_\delta$
exactly.
\end{definition}

\begin{remark}[Relation to~\eqref{eq:BS}]\label{rem:denom-vs-BS}
Lemma~\ref{lem:bohr-sommerfeld-pswf} gives
$\int\wncl\,dx = \pi(n+\tfrac{1}{2}) + O(1/n)$,
and Lemma~\ref{lem:denom-bs-pswf} gives
$\int 1/\wncl\,dx = \pi T/(n+\tfrac{1}{2}) + O(1/n^2)$.
The two are related by the discrete Bohr--Sommerfeld increment
argument of Lemma~\ref{lem:chi-increment}
(see Lemma~\ref{lem:denom-bs-pswf});
they are \emph{not} immediately equivalent by a symmetry argument.
\end{remark}

\subsection{Bohr--Sommerfeld quantization for $\wncl$}

\begin{lemma}[Bohr--Sommerfeld for PSWFs]%
\label{lem:bohr-sommerfeld-pswf}
For $n \le \gamma N$:
\begin{equation}\label{eq:BS}
  S^+_n := \int_{-x_+(n)}^{x_+(n)} \wncl(x)\,dx
  = \pi\bigl(n + \tfrac{1}{2}\bigr)
    + O\!\left(\frac{1}{n}\right).
\end{equation}
The error term $O(1/n)$ is uniform in $n \le \gamma N$,
with an implicit constant depending only on
$\delta$, $\gamma$, $c$, $T$.
Under stronger regularity assumptions on $\chi_n(c)$
the error can be sharpened to $O(1/n^2)$;
see \cite{OsipovRokhlinXiao}, Ch.~4, Thm.~4.7.

\emph{External input:} The WKB quantization result
\eqref{eq:BS} is imported from
\cite{OsipovRokhlinXiao}, Ch.~4, Thm.~4.7,
which is the authoritative reference for PSWF semiclassical
asymptotics. All subsequent results in this paper use only
the stated $O(1/n)$ bound as a black-box input;
no microlocal analysis, Airy-function matching, or
Riemann--Hilbert methods are reproduced here.
\end{lemma}

\begin{proof}
The PSWF eigenfunction $\psin$ satisfies the Sturm--Liouville
equation $(p\psin')' + \wncl^2/p\cdot\psin = 0$ on $I_\delta$.
Reducing to Liouville normal form \cite{Levitan}, Ch.~2,
produces a Schr\"odinger equation $u'' + [\wncl^2 - V]u = 0$
with potential $V = O(1)$ and $C^2$-norm $O(n^{-1})\cdot n = O(1)$
(using $\wncl \asymp n$, $(\wncl)' = O(1)$).
The standard WKB quantization
(\cite{OsipovRokhlinXiao}, Ch.~4, Thm.~4.7)
then gives \eqref{eq:BS} with error $O(1/n)$,
uniform in $n \le \gamma N$ since
$\wncl \ge c_{\delta,\gamma}n$ (Lemma~\ref{lem:G-regularity}(i))
and $\|\wncl\|_{C^2(I_\delta)} = O(n)$.
\end{proof}

\subsection{Discrete Bohr--Sommerfeld increment}

\begin{lemma}[Discrete Bohr--Sommerfeld increment for $\chi_n$]
\label{lem:chi-increment}
Let
\[
  S^+(\chi) := \int_{-x_+(\chi)}^{x_+(\chi)} \wncl(x;\chi)\,dx,
\]
which is $C^1$ and strictly increasing in $\chi$ on the bulk
regime (Lemma~\ref{lem:bohr-sommerfeld-pswf}).
Then for all $n \le \gamma N$,
\begin{align}
  S^+(\chi_{n+1}) - S^+(\chi_n)
    &= \pi + O\!\left(\tfrac{1}{n^2}\right),
  \label{eq:chi-increment}\\
  \chi_{n+1} - \chi_n
    &= 2n + O(1).
  \label{eq:chi-spacing}
\end{align}
\end{lemma}

\begin{proof}
\textbf{Step~1 (Bohr--Sommerfeld quantization).}
The eigenvalues $\chi_n$ satisfy
\[
  S^+(\chi_n) \;=\; \pi\!\left(n + \tfrac{1}{2}\right) + R(n),
  \qquad R(n) = O\!\left(\tfrac{1}{n}\right).
\]

\textbf{Step~2 (Remainder differencing).}
The remainder $R(n)$ has an expansion to first order in $1/n$,
\[
  R(n) \;=\; -\frac{\phi_0}{n} + O\!\left(\frac{1}{n^2}\right),
\]
with $\phi_0 = \phi_0(c,T) > 0$
(\cite{OsipovRokhlinXiao}, Ch.~4,
quantization refinement for bulk eigenvalues).
The error bound $O(1/n^2)$ holds uniformly in $n \le \gamma N$
and is stable under the unit shift $n \mapsto n+1$, since all
remainder estimates in the bulk regime are uniform over this
index range with constants depending only on
$(\delta, \gamma, c, T)$.
Hence by direct algebraic differencing,
\[
  R(n+1) - R(n)
  \;=\; \frac{\phi_0}{n(n+1)} + O\!\left(\frac{1}{n^3}\right)
  \;=\; O\!\left(\frac{1}{n^2}\right),
\]
which establishes~\eqref{eq:chi-increment} without any smoothness
assumption on $n$.

\textbf{Step~3 (Eigenvalue spacing, self-contained).}
From the asymptotic expansion
\[
  \chi_n(c) = n(n+1) + (2n+1)\frac{c}{\pi} + O\!\left(\frac{1}{n}\right),
\]
we compute directly:
\begin{align}
  \chi_{n+1} - \chi_n
  &= (n+1)(n+2) - n(n+1)
     + \frac{c}{\pi}\bigl((2n+3)-(2n+1)\bigr)
     + O(1) \notag \\
  &= 2n + 2 + \frac{2c}{\pi} + O(1)
   = 2n + O(1),
\end{align}
which establishes~\eqref{eq:chi-spacing}.
\end{proof}

%%------------------------------------------------------------------
%% NEU: Stability of inverse Bohr--Sommerfeld map
%%------------------------------------------------------------------
\begin{lemma}[Stability of inverse Bohr--Sommerfeld map]
\label{lem:bs-inverse-stability}
The map $\chi \mapsto S^+(\chi)$ is $C^1$ and strictly
increasing on the bulk regime, with
\begin{equation}\label{eq:dSplus-lower}
  \frac{dS^+}{d\chi}(\chi) \;=\; \frac{S^-(\chi)}{2T^2}
  \;\ge\; \frac{c_{\delta,\gamma}}{n},
  \qquad \chi \in [\chi_n, \chi_{n+1}],
\end{equation}
uniformly in $n \le \gamma N$.
Consequently, the inverse map $S^+ \mapsto \chi$ satisfies
\begin{equation}\label{eq:chi-from-Splus}
  \chi_{n+1} - \chi_n
  = \frac{S^+(\chi_{n+1}) - S^+(\chi_n)}{(dS^+/d\chi)(\chi_n^*)}
  = \frac{\pi + O(1/n^2)}{S^-(\chi_n^*)/(2T^2)}
\end{equation}
for some $\chi_n^* \in (\chi_n, \chi_{n+1})$ (MVT),
and the $O(1/n^2)$ error in~\eqref{eq:chi-increment} is stable
under evaluation at $\chi_n^*$ versus $\chi_n$, in the sense that
\begin{equation}\label{eq:Sminus-stability}
  \left|\frac{dS^-}{d\chi}(\chi)\right|
  \;\le\; \frac{C_{\delta,\gamma}}{n^2},
  \qquad \chi \in [\chi_n, \chi_{n+1}],
\end{equation}
so that $S^-(\chi_n^*) = S^-(\chi_n) + O(1/n)$,
and the resulting error in $\chi_{n+1} - \chi_n$
is $O(1)$ as stated in~\eqref{eq:chi-spacing}.

In particular, all MVT applications in the proofs of
Lemmas~\ref{lem:chi-increment} and~\ref{lem:denom-bs-pswf}
are applied only to the \emph{continuous} function
$\chi \mapsto S^+(\chi)$, never to $n \mapsto \chi_n$
as a continuous function of $n$.
No smoothness of $\chi_n$ in $n$ is assumed or used.
\end{lemma}

\begin{proof}
\textbf{Step~1 ($dS^+/d\chi$ formula).}
Differentiating under the integral (boundary terms vanish at
turning points where $\wncl = 0$):
\[
  \frac{dS^+}{d\chi}(\chi)
  = \int_{-x_+}^{x_+} \frac{\partial \wncl}{\partial \chi}\,dx
  = \int_{-x_+}^{x_+} \frac{1}{2T^2\,\wncl(x;\chi)}\,dx
  = \frac{S^-(\chi)}{2T^2}.
\]

\textbf{Step~2 (Lower bound).}
By Lemma~\ref{lem:G-regularity}(i), $\wncl \le C_{\delta,\gamma}n$
on $I_\delta$, so $1/\wncl \ge c_{\delta,\gamma}/n$, giving
$S^-(\chi) \ge 2|I_\delta|\cdot c_{\delta,\gamma}/n$,
hence $dS^+/d\chi \ge c_{\delta,\gamma}/n > 0$.

\textbf{Step~3 ($dS^-/d\chi$ bound).}
Differentiating $S^-(\chi) = \int_{-x_+}^{x_+} 1/\wncl\,dx$:
\[
  \frac{dS^-}{d\chi}(\chi)
  = -\int_{-x_+}^{x_+} \frac{1}{2T^2\,\wncl^3}\,dx.
\]
Since $\wncl \ge c_{\delta,\gamma}n$ on $I_\delta$:
$|dS^-/d\chi| \le C_{\delta,\gamma}/(n^3)\cdot 2T
= O(1/n^3)$,
which is $O(1/n^2)$ as stated (the bound is in fact sharper).

\textbf{Step~4 (MVT stability).}
Since $|\chi_{n+1} - \chi_n| = O(n)$
(Lemma~\ref{lem:chi-increment}, \eqref{eq:chi-spacing}):
\[
  |S^-(\chi_n^*) - S^-(\chi_n)|
  \le \frac{C_{\delta,\gamma}}{n^3}\cdot O(n)
  = O\!\left(\frac{1}{n^2}\right).
\]
This is absorbed into the $O(1/n)$ precision of $S^-_n$
from Lemma~\ref{lem:denom-bs-pswf},
so the evaluation point $\chi_n^*$ vs.\ $\chi_n$
introduces no additional error at leading order.
\end{proof}

\subsection{Bohr--Sommerfeld quantization for $1/\wncl$}

\begin{lemma}[Dual Bohr--Sommerfeld: $\int 1/\wncl$]%
\label{lem:denom-bs-pswf}
For $n \le \gamma N$:
\begin{equation}\label{eq:denom-BS-lemma}
  S_n^- := \int_{-x_+(n)}^{x_+(n)} \frac{dx}{\wncl(x)}
  = \frac{\pi T}{n + \tfrac{1}{2}}
    + O\!\left(\frac{1}{n^2}\right),
\end{equation}
uniformly in $n \le \gamma N$.
\end{lemma}

\begin{proof}
Recall from \eqref{eq:wncl-def} that
$\wncl(x) = \wncl(x;\chi_n)$
depends on the eigenvalue parameter $\chi = \chi_n(c)$.
Differentiate $S^+(\chi)$ with respect to $\chi$:
\begin{equation}\label{eq:diff-chi}
  \frac{d}{d\chi} S^+(\chi)
  = \frac{d}{d\chi}\int_{-x_+}^{x_+} \wncl(x;\chi)\,dx.
\end{equation}
The boundary terms from differentiating the limits cancel:
at $x = \pm x_+$, $\wncl(x_+(\chi);\chi) = 0$ by definition
of the turning point, so the endpoint contributions vanish.
Differentiating under the integral sign (valid by dominated
convergence since $\partial\wncl/\partial\chi
= 1/(2T^2\wncl) \in L^1$ on the interior):
\begin{equation}\label{eq:dSplus-dchi}
  \frac{d}{d\chi} S^+(\chi)
  = \int_{-x_+}^{x_+} \frac{\partial\wncl}{\partial\chi}\,dx
  = \int_{-x_+}^{x_+} \frac{1}{2T^2\,\wncl(x;\chi)}\,dx
  = \frac{1}{2T^2}\,S^-(\chi).
\end{equation}
From Lemma~\ref{lem:chi-increment} (discrete BS increment):
\[
  S^+(\chi_{n+1}) - S^+(\chi_n) = \pi + O\!\left(\frac{1}{n^2}\right),
\]
and there exists $\chi_n^* \in (\chi_n, \chi_{n+1})$ with
\[
  \frac{dS^+}{d\chi}(\chi_n^*) \cdot (\chi_{n+1} - \chi_n)
  = \pi + O\!\left(\frac{1}{n^2}\right).
\]
By Lemma~\ref{lem:bs-inverse-stability}~\eqref{eq:Sminus-stability},
the evaluation at $\chi_n^*$ versus $\chi_n$ introduces only
an $O(1/n^2)$ error in $S^-$, so substituting
\eqref{eq:dSplus-dchi} at $\chi = \chi_n^*$:
\[
  \frac{S^-(\chi_n)}{2T^2} \cdot (\chi_{n+1} - \chi_n)
  = \pi + O\!\left(\frac{1}{n^2}\right).
\]
Using $\chi_{n+1} - \chi_n = 2n + O(1)$
(Lemma~\ref{lem:chi-increment}, \eqref{eq:chi-spacing}):
\[
  \frac{S_n^-}{2T^2}\cdot(2n + O(1))
  = \pi + O\!\left(\frac{1}{n^2}\right).
\]

\textit{Normalization correction.}
Set $t = x/T$. Then $dx = T\,dt$ and
$\wncl(x) = \wncl(Tt)$.
Hence
\[
  \int_{-x_+(n)}^{x_+(n)} \frac{dx}{\wncl(x)}
  = T \int_{-t_+(n)}^{t_+(n)} \frac{dt}{\wncl(Tt)}.
\]
Applying Lemma~\ref{lem:bohr-sommerfeld-pswf} in the
normalized variable and using
$2n + O(1) = 2(n+\tfrac{1}{2}) + O(1/n)$:
\begin{equation}\label{eq:Sminus-final}
  S_n^- := \int_{-x_+(n)}^{x_+(n)} \frac{dx}{\wncl(x)}
  = \frac{\pi T}{n + \tfrac{1}{2}}
    + O\!\left(\frac{1}{n^2}\right).
\end{equation}
This is \eqref{eq:denom-BS-lemma}. No hidden $T^2$-factors
remain: the factor $T$ on the right-hand side reflects
the correct dimensional analysis for $x \in [-T,T]$,
consistent with WKB conventions for the non-normalized variable.
\end{proof}

\begin{remark}[Self-containedness]
The proof of Lemma~\ref{lem:denom-bs-pswf} uses only
Lemma~\ref{lem:bohr-sommerfeld-pswf} (\eqref{eq:BS}),
the discrete increment Lemma~\ref{lem:chi-increment},
and Lemma~\ref{lem:bs-inverse-stability} for MVT stability.
In particular, no smoothness of $\chi_n$ as a function of $n$
is assumed: the MVT is applied only to the $C^1$ function
$\chi \mapsto S^+(\chi)$, and the evaluation-point shift
$\chi_n^* \to \chi_n$ is controlled explicitly by
\eqref{eq:Sminus-stability}.
It does not invoke microlocal analysis, Airy functions,
or Riemann--Hilbert methods.
\end{remark}

\subsection{Regularity of the WKB coefficient}

\begin{lemma}[Regularity of $G$ and $\wncl$]%
\label{lem:G-regularity}
Define:
\begin{equation}\label{eq:G-def}
  G(x) := \frac{(\wncl)'(x)}{2\wncl(x)} + \frac{p'(x)}{4p(x)},
  \qquad p(x) := 1 - x^2/T^2.
\end{equation}
Then on $I_\delta$, uniformly in $n \le \gamma N$:
\begin{enumerate}[(i)]
  \item $\wncl(x) \ge c_{\delta,\gamma}\,n$,
  \item $|(\wncl)'(x)| \le C_{\delta,\gamma}$,
  \item $\|G\|_{L^\infty(I_\delta)} \le C_{\delta,\gamma}/n$,
  \item $\|G'\|_{L^1(I_\delta)} \le C_{\delta,\gamma}/n$,
  \item $p(x) \ge \delta^2$ and
        $|p'(x)| = 2|x|/T^2 \le 2/T$.
\end{enumerate}
\end{lemma}

\begin{proof}
(i): $\wncl(x)^2 \ge \chi_n/T^2 \ge C_\gamma n^2/T^2$,
using $\chi_n \ge C_\gamma n^2$ for $n \le \gamma N$.
(ii): Direct differentiation of \eqref{eq:wncl-def}.
(iii)--(iv): $G = (\wncl)'/(2\wncl) + p'/(4p)$;
each term is $O(1/n)$ on $I_\delta$ by (i),(ii),(v).
(v): Explicit from $p(x) = 1 - x^2/T^2 \ge \delta^2$
on $I_\delta$.
\end{proof}

\subsection{Turning point cutoff}

\begin{lemma}[Exponential turning point cutoff]%
\label{lem:turning-point-cutoff}
For $n \le \gamma N$ and $f \in L^\infty([-T,T])$:
\begin{equation}\label{eq:cutoff}
  \left|\int_{[-T,T]\setminus I_\delta}
    f(x)\,\psin(x)^2\,dx\right|
  \le C_{\delta,\gamma}\,e^{-\beta_{\delta,\gamma} n}\,
     \lambda_n\,\|f\|_{L^\infty},
\end{equation}
where $\beta_{\delta,\gamma} > 0$ depends only on
$\delta$, $\gamma$, $c$, $T$.
\end{lemma}

\begin{proof}
$\psin$ decays exponentially outside $[-x_+(n), x_+(n)]$,
and since $x_+(n) \le (1-\eta_\gamma)T < (1-\delta)T$,
the region $[-T,T]\setminus I_\delta$ lies strictly outside
the turning points; see \cite{OsipovRokhlinXiao}, Ch.~5.
\end{proof}

%%=================================================================
\section{Pr\"ufer Transform}\label{sec:prufer}
%%=================================================================

\subsection{Pr\"ufer substitution}

On $I_\delta$, $\psin$ satisfies:
\begin{equation}\label{eq:SL}
  (p(x)\psin')' + \wncl(x)^2/p(x)\cdot\psin = 0.
\end{equation}
Define the Pr\"ufer variables $(r_n, \theta_n)$ by:
\begin{equation}\label{eq:prufer-def}
  \psin(x) = r_n(x)\sin(\theta_n(x)),
  \quad
  p(x)\psin'(x) = r_n(x)\wncl(x)\cos(\theta_n(x)).
\end{equation}

\subsection{Pr\"ufer ODEs}

\begin{lemma}[Pr\"ufer system]\label{lem:prufer-ode}
The Pr\"ufer variables satisfy on $I_\delta$:
\begin{align}
  \theta_n'
  &= \wncl(x)/p(x) + G(x)\sin(2\theta_n(x)),
  \label{eq:theta-ode}\\
  (\log r_n)'
  &= G(x)\cos(2\theta_n(x)) + D(x),
  \label{eq:r-ode}
\end{align}
where $D(x) := p'(x)/(4p(x))$ and $G(x)$ is as in
Lemma~\ref{lem:G-regularity}.
\end{lemma}

\begin{proof}
Standard Pr\"ufer differentiation; see \cite{Levitan}, Ch.~1.
\end{proof}

\subsection{Phase lower bound}

\begin{lemma}[Phase lower bound]\label{lem:phase-lower}
On $I_\delta$, uniformly in $n \le \gamma N$:
\begin{equation}\label{eq:phase-lb}
  \theta_n'(x) \ge \frac{c_{\delta,\gamma}\,n}{T}.
\end{equation}
\end{lemma}

\begin{proof}
From \eqref{eq:theta-ode} and
Lemma~\ref{lem:G-regularity}(i),(iii):
$\theta_n' \ge \wncl/p - |G|
\ge c_{\delta,\gamma}n - C_{\delta,\gamma}/n
\ge c_{\delta,\gamma}n/2$
for $n$ sufficiently large;
adjusting the constant absorbs the finite cases.
\end{proof}

\subsection{Phase upper bound}

\begin{lemma}[Phase upper bound]\label{lem:phase-upper}
On $I_\delta$, uniformly in $n \le \gamma N$:
\begin{equation}\label{eq:phase-ub}
  \theta_n'(x) \le \frac{C_{\delta,\gamma}\,n}{T}.
\end{equation}
Consequently $\theta_n'(x) \asymp n$ uniformly on $I_\delta$.
\end{lemma}

\begin{proof}
From \eqref{eq:theta-ode}:
$\theta_n' = \wncl(x)/p(x) + G(x)\sin(2\theta_n(x))$.
By Lemma~\ref{lem:G-regularity}(i): $\wncl(x) \le C_{\delta,\gamma}n$.
By Lemma~\ref{lem:G-regularity}(v): $p(x) \ge \delta^2 > 0$.
By Lemma~\ref{lem:G-regularity}(iii): $|G| \le C_{\delta,\gamma}/n$.
Hence:
\[
  \theta_n'(x)
  \le \frac{C_{\delta,\gamma}n}{\delta^2}
     + \frac{C_{\delta,\gamma}}{n}
  \le \frac{C_{\delta,\gamma}'\,n}{T}.
\]
\end{proof}

\begin{remark}[Phase quasi-linearity]
Lemmas~\ref{lem:phase-lower} and~\ref{lem:phase-upper} together
give $c_{\delta,\gamma}n \le T\theta_n'(x) \le C_{\delta,\gamma}n$,
i.e., $\theta_n$ is \emph{quasi-linear} with rate $\asymp n$.
This is used:
(a) in Lemma~\ref{lem:prufer-oscillation-control}
    (lower bound, for the IBP denominator);
(b) in Lemma~\ref{lem:pointwise-control}, Step~2
    (upper bound, for $|J| \ge c_{\delta,\gamma}'/n$);
(c) in Lemma~\ref{lem:phase-second-deriv}
    (ratio $\theta_n''/(\theta_n')^2 = O(1/n^2)$).
\end{remark}

\subsection{Phase second derivative bound}

\begin{lemma}[Phase second derivative]%
\label{lem:phase-second-deriv}
On $I_\delta$, uniformly in $n \le \gamma N$:
\begin{equation}\label{eq:phase-2nd}
  |\theta_n''(x)| \le C_{\delta,\gamma},
  \qquad
  \left|\frac{\theta_n''(x)}{(\theta_n'(x))^2}\right|
  \le \frac{C_{\delta,\gamma}}{n^2}.
\end{equation}
\end{lemma}

\begin{proof}
Differentiate \eqref{eq:theta-ode}:
$\theta_n'' = (\wncl/p)' + (G\sin(2\theta_n))'$.
By Lemma~\ref{lem:G-regularity},
$(\wncl/p)' = O(1)$, $G' = O(1/n)$,
$G\theta_n' = O(1)$, so $|\theta_n''| \le C_{\delta,\gamma}$.
The ratio bound follows from
$\theta_n' \ge c_{\delta,\gamma}n$
(Lemma~\ref{lem:phase-lower}).
\end{proof}

%%=================================================================
\section{Oscillation Control Lemma}\label{sec:oscillation}
%%=================================================================

\begin{lemma}[Pr\"ufer oscillation control]%
\label{lem:prufer-oscillation-control}
Let $h \in C^1(I_\delta)$. Then for every
$x_0, x \in I_\delta$:
\begin{equation}\label{eq:osc-control-sup}
  \sup_{x_0,\,x \in I_\delta}
  \left|\int_{x_0}^{x} h(t)\cos(2\theta_n(t))\,dt\right|
  \le \frac{C_{\delta,\gamma}}{n}
      \left(\|h\|_{L^\infty(I_\delta)}
        + \|h'\|_{L^1(I_\delta)}\right).
\end{equation}
In particular:
\begin{equation}\label{eq:osc-control}
  \left|\int_{I_\delta} h(x)\cos(2\theta_n(x))\,dx\right|
  \le \frac{C_{\delta,\gamma}}{n}
      \left(\|h\|_{L^\infty(I_\delta)}
        + \|h'\|_{L^1(I_\delta)}\right).
\end{equation}
\end{lemma}

\begin{proof}
Fix $x_0 < x$ in $I_\delta$.
Since $\theta_n' \ge (c_{\delta,\gamma}/T)n > 0$
(Lemma~\ref{lem:phase-lower}),
integrate by parts:
\begin{equation}\label{eq:ibp2}
  \int_{x_0}^{x} h\cos(2\theta_n)\,dt
  = \left[\frac{h(t)\sin(2\theta_n(t))}{2\theta_n'(t)}\right]_{x_0}^{x}
    - \int_{x_0}^{x} \sin(2\theta_n)\cdot
      \frac{d}{dt}\!\left(\frac{h}{2\theta_n'}\right)\,dt.
\end{equation}
\textit{Boundary term:}
$|\cdot| \le T\|h\|_{L^\infty}/(c_{\delta,\gamma}n)$,
uniform in $x_0, x$.

\textit{$h'$-term:}
$|\cdot| \le \frac{T}{2c_{\delta,\gamma}n}\|h'\|_{L^1(I_\delta)}$.

\textit{$\theta_n''$-term:}
By Lemma~\ref{lem:phase-second-deriv},
$|\theta_n''|/(\theta_n')^2 \le C_{\delta,\gamma}/n^2$, so
$|\cdot| \le 2TC_{\delta,\gamma}\|h\|_{L^\infty}/n^2$.

All three bounds are uniform in $x_0, x \in I_\delta$,
giving \eqref{eq:osc-control-sup}.
\end{proof}

\begin{remark}[Uniformity in $x_0, x$ is essential]
The sup-form \eqref{eq:osc-control-sup} is required in
Lemma~\ref{lem:amplitude-drift}, Step~3,
where one needs
$\sup_{x}|\int_{x_0}^x G\cos(2\theta_n)| \le C/n$
uniformly in $x \in I_\delta$.
\end{remark}

\begin{remark}[Sharpness]
For $h \equiv 1$ and $\theta_n(x) \approx nx/T$,
$|\int_{x_0}^x\cos(2nt/T)\,dt|
= O(T/n)$,
so \eqref{eq:osc-control-sup} is sharp in $n$.
\end{remark}

\begin{corollary}[Oscillation control with slowly varying weight]%
\label{cor:osc-slowly-varying}
If $a \in C^1(I_\delta)$ and $b_n \in C^1(I_\delta)$ satisfies
$\|b_n\|_{L^\infty} \le C$ and
$\|b_n'\|_{L^\infty} \le C/n$,
uniformly in $n \le \gamma N$, then:
\begin{equation}\label{eq:osc-sv}
  \left|\int_{I_\delta}
    a(x)\,b_n(x)\,\cos(2\theta_n(x))\,dx\right|
  = O\!\left(\frac{\|a\|_{C^1(I_\delta)}\,
                   \|b_n\|_{L^\infty}}{n}\right).
\end{equation}
\end{corollary}

\begin{proof}
Set $h := ab_n$.
Then $\|h\|_{L^\infty} \le \|a\|_{L^\infty}\|b_n\|_{L^\infty}$
and $\|h'\|_{L^1} \le 2T(\|a'\|_{L^\infty}\|b_n\|_{L^\infty}
+ \|a\|_{L^\infty}\cdot C/n)
= O(\|a\|_{C^1}\|b_n\|_{L^\infty})$.
Apply Lemma~\ref{lem:prufer-oscillation-control}.
\end{proof}

\begin{remark}[Integration of oscillatory control into amplitude dynamics]
\label{rem:osc-gronwall-link}
Lemma~\ref{lem:prufer-oscillation-control} provides a uniform
bound of the form
\[
  \left|
    \int_{x_0}^{x} G(t)\cos(2\theta_n(t))\,dt
  \right|
  \le \frac{C_{\delta,\gamma}}{n},
  \qquad \forall\, x \in I_\delta.
\]
In the amplitude deviation equation~\eqref{eq:deviation-ode},
this estimate controls the oscillatory forcing term in the sense
that
\[
  \left|
    \int_{x_0}^{x} \widetilde{B}_n(t)\,dt
  \right|
  \le \frac{C_{\delta,\gamma}}{n}.
\]
Consequently, the amplitude equation fits into a Gronwall-type
stability framework with forcing of size $O(1/n)$ and a
coefficient whose integral over $I_\delta$ is $O(1)$.
This ensures that the oscillatory cancellation established in
Section~\ref{sec:oscillation} propagates directly into the
Gronwall estimate in Lemma~\ref{lem:amplitude-drift},
yielding a uniform $O(1/n)$ control of the amplitude deviation.
\end{remark}

%%=================================================================
\section{Amplitude Control}\label{sec:amplitude}
%%=================================================================

\begin{definition}[WKB amplitude]\label{def:wkb-amp}
Fix $x_0 \in I_\delta$. Define:
\begin{equation}\label{eq:wkb-amp}
  r_n^{\mathrm{WKB}}(x)
  := r_n(x_0)\sqrt{\frac{\wncl(x_0)}{\wncl(x)}},
  \quad x \in I_\delta,
\end{equation}
the solution of the averaged Pr\"ufer amplitude ODE
$(\log r)' = D(x) - (\wncl)'/(2\wncl)$
with $r(x_0) = r_n(x_0)$.
Note:
\begin{equation}\label{eq:wkb-amp-deriv}
  \frac{d}{dx}\bigl[r_n^{\mathrm{WKB}}(x)^2\bigr]
  = O\!\left(\frac{r_n(x_0)^2}{n}\right),
\end{equation}
uniformly on $I_\delta$,
since $(\wncl)' = O(1)$ and $\wncl \asymp n$
(Lemma~\ref{lem:G-regularity}(i),(ii)).
\end{definition}

\begin{lemma}[Exact amplitude deviation ODE]
\label{lem:deviation-ode}
Let $\Delta_n(x) := r_n(x)^2 - r_n^{\mathrm{WKB}}(x)^2$.
Then:
\begin{equation}\label{eq:deviation-ode}
  \Delta_n'(x)
  = 2D(x)\,\Delta_n(x)
    + 2r_n^{\mathrm{WKB}}(x)^2\,\widetilde{B}_n(x)
    + 2\Delta_n(x)\,G(x)\cos(2\theta_n(x)),
\end{equation}
where:
\begin{equation}\label{eq:Btilde-def}
  \widetilde{B}_n(x)
  := \frac{(\wncl)'(x)}{2\wncl(x)}\cos(2\theta_n(x))
     + G(x)\cos(2\theta_n(x)),
\end{equation}
with $|\widetilde{B}_n(x)| \le C_{\delta,\gamma}/n$
uniformly on $I_\delta$.
The $p'/(4p)$-contributions cancel exactly as a consequence
of the definition of $r_n^{\mathrm{WKB}}$,
leaving oscillatory forcing $O(1/n)$ and a linear term
$2\Delta_n G\cos(2\theta_n)$ with coefficient $O(1/n)$.
\end{lemma}

\begin{proof}
$(r_n^2)' = 2r_n^2(D + G\cos(2\theta_n))$
from \eqref{eq:r-ode}.
$(r_n^{\mathrm{WKB},2})' = 2r_n^{\mathrm{WKB},2}(D - (\wncl)'/(2\wncl))$
from Definition~\ref{def:wkb-amp}.
Subtracting and collecting terms gives \eqref{eq:deviation-ode}
with $\widetilde{B}_n$ as in \eqref{eq:Btilde-def}.
The $D$-terms cancel by the definition of $r_n^{\mathrm{WKB}}$.
\end{proof}

\begin{lemma}[Pointwise amplitude control]%
\label{lem:pointwise-control}
For any fixed $x_0 \in I_\delta$ and all $x \in I_\delta$,
uniformly in $n \le \gamma N$:
\begin{equation}\label{eq:amp-comp}
  r_n(x)^2 \le C_{\delta,\gamma}\,r_n(x_0)^2.
\end{equation}
Consequently:
\begin{equation}\label{eq:pointwise-lambda}
  r_n(x_0)^2 \le C_{\delta,\gamma}\,\lambda_n.
\end{equation}
\end{lemma}

\begin{proof}
\textit{Step~1: Amplitude comparability.}
From \eqref{eq:r-ode}:
$|(\log r_n)'| \le |G| + |D|
\le C_{\delta,\gamma}/n + C_{\delta,\gamma}
\le C_{\delta,\gamma}$,
uniformly on $I_\delta$.
Integrating over $I_\delta$ with $|I_\delta| \le 2T$:
\[
  \left|\log\frac{r_n(x)}{r_n(x_0)}\right|
  \le 2C_{\delta,\gamma}T
  =: K_{\delta,\gamma},
\]
so $r_n(x)^2 \le e^{2K_{\delta,\gamma}}r_n(x_0)^2
=: C_{\delta,\gamma}\,r_n(x_0)^2$.

\textit{Step~2: Global $L^2$ lower bound via explicit $J$.}
Since $\theta_n' \ge (c_{\delta,\gamma}/T)n$
(Lemma~\ref{lem:phase-lower}),
$\theta_n$ traverses at least
$M := \lfloor c_{\delta,\gamma}n/\pi\rfloor \ge 1$
complete periods on $I_\delta$ for $n \ge n_0(\delta,\gamma)$.

Fix any complete period $[a,b] \subset I_\delta$
with $\theta_n(b) - \theta_n(a) = \pi$.
The subinterval
\[
  J := \theta_n^{-1}\!\left(
         \theta_n(a) + \tfrac{\pi}{6},\;
         \theta_n(a) + \tfrac{5\pi}{6}
       \right) \subset [a,b]
\]
satisfies $\sin^2\theta_n(x) \ge 1/4$ for all $x \in J$,
and by Lemma~\ref{lem:phase-upper}:
\[
  |J| \ge \frac{2\pi/3}{C_{\delta,\gamma}n}
      =: \frac{c_{\delta,\gamma}'}{n}.
\]
By \eqref{eq:amp-comp}:
$\psin^2 = r_n^2\sin^2\theta_n
\ge r_n^2/4 \ge r_n(x_0)^2/(4C_{\delta,\gamma})$
on $J$, so:
\[
  \lambda_n
  \ge \int_J\psin^2
  \ge \frac{c_{\delta,\gamma}'}{n}
     \cdot\frac{r_n(x_0)^2}{4C_{\delta,\gamma}},
\]
giving \eqref{eq:pointwise-lambda}.
\end{proof}

\begin{remark}[Non-circularity of Lemma~\ref{lem:pointwise-control}]
This lemma uses only the Pr\"ufer ODE~\eqref{eq:r-ode},
Lemmas~\ref{lem:G-regularity}, \ref{lem:phase-lower},
and~\ref{lem:phase-upper},
and the $L^2$-normalization $\int\psin^2 = \lambda_n$.
It does not use Lemmas~\ref{lem:amplitude-drift}
or~\ref{lem:normalization},
nor Theorem~\ref{thm:weak-limit}.
\end{remark}

\begin{lemma}[Amplitude drift]\label{lem:amplitude-drift}
For $n \le \gamma N$ and all $x \in I_\delta$:
\begin{align}
  \left|\frac{r_n(x)^2}{r_n^{\mathrm{WKB}}(x)^2} - 1\right|
  &\le \frac{C_{\delta,\gamma}}{n},
  \label{eq:amp-ratio}\\
  \left|\Delta_n'(x)\right|
  &\le \frac{C_{\delta,\gamma}\,r_n(x_0)^2}{n}.
  \label{eq:amp-dev-deriv}
\end{align}
\end{lemma}

\begin{remark}[Two bounds, two roles]\label{rem:two-bounds}
Bound~\eqref{eq:amp-ratio} gives pointwise ratio control,
used in Lemmas~\ref{lem:normalization}
and~\ref{thm:weak-limit}.
Bound~\eqref{eq:amp-dev-deriv} shows
$\|b_n'\|_{L^\infty} = O(1/n)$
for $b_n := r_n^2/r_n(x_0)^2$,
as required by Corollary~\ref{cor:osc-slowly-varying}.
\end{remark}

\begin{proof}
\textit{Step~1: Drift cancellation identity.}
From \eqref{eq:r-ode} and Definition~\ref{def:wkb-amp}:
\[
  \frac{d}{dx}\log\frac{r_n(x)}{r_n^{\mathrm{WKB}}(x)}
  = G(x)\bigl(1 + \cos(2\theta_n(x))\bigr).
\]
The $D$-terms cancel exactly by definition of $r_n^{\mathrm{WKB}}$.

\textit{Step~2: Reduction to oscillatory integral.}
Integrating from $x_0$ and splitting
$G(1+\cos(2\theta_n)) = G + G\cos(2\theta_n)$:
\begin{equation}\label{eq:drift-oscillatory}
  \log\frac{r_n(x)}{r_n^{\mathrm{WKB}}(x)}
  = \int_{x_0}^x G(t)\cos(2\theta_n(t))\,dt,
\end{equation}
where the $\int G\,dt$ term cancels against the WKB definition.

\textit{Step~3: Proof of \eqref{eq:amp-ratio}.}
By Lemma~\ref{lem:prufer-oscillation-control}
(sup-form, $h = G$,
$\|G\|_{L^\infty} \le C_{\delta,\gamma}/n$,
$\|G'\|_{L^1} \le C_{\delta,\gamma}/n$):
\[
  \sup_{x \in I_\delta}
  \left|\int_{x_0}^x G(t)\cos(2\theta_n(t))\,dt\right|
  \le \frac{C_{\delta,\gamma}}{n}.
\]
Exponentiating and squaring \eqref{eq:drift-oscillatory}
gives \eqref{eq:amp-ratio}.

\textit{Step~4: Proof of \eqref{eq:amp-dev-deriv} via Gronwall.}
By Lemma~\ref{lem:deviation-ode}, $\Delta_n$ satisfies
\eqref{eq:deviation-ode} with $\Delta_n(x_0) = 0$,
forcing $|F| \le C_{\delta,\gamma}\,r_n(x_0)^2/n$,
and coefficient $A(x) := 2D(x) + 2G(x)\cos(2\theta_n(x))$
with $\int_{I_\delta} A\,dx = O(1)$ (the $D$-part contributes
$O(1)$, the oscillatory $G$-part contributes $O(1/n)$ by
Lemma~\ref{lem:prufer-oscillation-control};
see Remark~\ref{rem:osc-gronwall-link}).
The Gronwall--Bellman inequality \cite{Bellman1943} gives:
\[
  |\Delta_n(x)|
  \le \frac{C_{\delta,\gamma}'\,r_n(x_0)^2}{n},
\]
and substituting back into \eqref{eq:deviation-ode}
establishes \eqref{eq:amp-dev-deriv}.
\end{proof}

%%=================================================================
\section{Normalization Matching}\label{sec:normalization}
%%=================================================================

\begin{lemma}[WKB normalization]\label{lem:normalization}
For $n \le \gamma N$:
\begin{align}
  \frac{1}{2}\int_{I_\delta} r_n^{\mathrm{WKB}}(x)^2\,dx
  &= \lambda_n\left(1 + O\!\left(\frac{1}{n}\right)\right),
  \label{eq:norm-match}\\
  \frac{r_n^{\mathrm{WKB}}(x)^2}{2}
  &= \lambda_n\,\rhocl_n(x)
     \left(1 + O\!\left(\frac{1}{n}\right)\right).
  \label{eq:rn-vs-rhocl}
\end{align}
\end{lemma}

\begin{proof}
\textit{Step~1: Global normalization.}
Using $\sin^2\theta_n = \tfrac{1}{2} - \tfrac{1}{2}\cos(2\theta_n)$:
\[
  \lambda_n
  = \frac{1}{2}\int_{I_\delta}r_n^2
    - \frac{1}{2}\int_{I_\delta}r_n^2\cos(2\theta_n)
    + O(e^{-\beta n}\lambda_n).
\]
Set $b_n := r_n^2/r_n(x_0)^2$.
By \eqref{eq:amp-ratio}: $\|b_n\|_{L^\infty} \le C$.
By Remark~\ref{rem:two-bounds}: $\|b_n'\|_{L^\infty} = O(1/n)$.
Corollary~\ref{cor:osc-slowly-varying} gives
$|\int_{I_\delta}r_n^2\cos(2\theta_n)|
\le C_{\delta,\gamma}\lambda_n/n$,
hence \eqref{eq:norm-match}.

\textit{Step~2: Pointwise ratio to $\rhocl_n$.}
From \eqref{eq:wkb-amp}:
$r_n^{\mathrm{WKB}}(x)^2 = r_n(x_0)^2\,\wncl(x_0)/\wncl(x)$.
Integrating over $I_\delta$ and using \eqref{eq:norm-match},
Lemma~\ref{lem:denom-bs-pswf}, and
Definition~\ref{def:rho-cl}:
\[
  \frac{r_n^{\mathrm{WKB}}(x)^2}{2}
  = \frac{\lambda_n(1+O(1/n))}
         {\wncl(x)\cdot S_n^-}
  = \lambda_n\,\rhocl_n(x)(1+O(1/n)).
\]
\end{proof}

%%=================================================================
\section{Proof of the Main Theorem}\label{sec:weaklimit}
%%=================================================================

\begin{theorem}[Semiclassical expectation asymptotics]%
\label{thm:weak-limit}
For $n \le \gamma N$ (in the regime~\eqref{eq:semiclassical-regime})
and every $f \in C^1([-T,T])$:
\begin{equation}\label{eq:weak-limit}
  \int_{-T}^T f(x)\,\psin(x)^2\,dx
  = \lambda_n(c)
    \int_{-x_+(n)}^{x_+(n)} f(x)\,\rhocl_n(x)\,dx
    + R_n(f),
\end{equation}
where:
\begin{equation}\label{eq:remainder}
  |R_n(f)|
  \le \frac{C_{\delta,\gamma}\,\lambda_n(c)}{n}
      \left(\|f\|_{L^\infty} + T\,\|f'\|_{L^\infty}\right).
\end{equation}
The constant $C_{\delta,\gamma}$ depends only on
$\delta$, $\gamma$, $c$, $T$, not on $n$ or $f$.
This is a \emph{uniform semiclassical matching principle}:
for each fixed $f \in C^1$, the spectral integral
$\int f \psin^2$ is approximated by
$\lambda_n \int f \rhocl_n$ with explicit rate $O(\lambda_n/n)$,
uniform over the bulk regime $n \le \gamma N$.
\end{theorem}

\begin{proof}
\textit{Step~1 (Bulk cutoff).}
By Lemma~\ref{lem:turning-point-cutoff}:
$\int_{-T}^T f\psin^2
= \int_{I_\delta} f\psin^2
+ O(e^{-\beta n}\lambda_n\|f\|_{L^\infty})$.

\textit{Step~2 (Pr\"ufer decomposition).}
$\int_{I_\delta}f\psin^2
= \tfrac{1}{2}\int_{I_\delta}fr_n^2
- \tfrac{1}{2}\int_{I_\delta}fr_n^2\cos(2\theta_n)$.

\textit{Step~3 (Oscillation term).}
Set $a = f$, $b_n = r_n^2/r_n(x_0)^2$.
Corollary~\ref{cor:osc-slowly-varying} gives:
\[
  \left|\int_{I_\delta}fr_n^2\cos(2\theta_n)\right|
  \le \frac{C_{\delta,\gamma}\lambda_n
            (\|f\|_{L^\infty} + T\|f'\|_{L^\infty})}{n}.
\]

\textit{Step~4 (Mean term).}
By \eqref{eq:rn-vs-rhocl}:
$\tfrac{1}{2}\int_{I_\delta}fr_n^2
= \lambda_n\int_{-x_+(n)}^{x_+(n)}f\rhocl_n
+ O(\lambda_n\|f\|_{L^\infty}/n)$.

\textit{Step~5 (Assembly).}
All exponential contributions satisfy $e^{-\beta n} \le C/n$
for $n \ge n_0$, so they are absorbed into $O(\lambda_n/n)$.
Combining Steps~1--5 gives
\eqref{eq:weak-limit}--\eqref{eq:remainder}.
\end{proof}

\begin{remark}[Uniformity of the constant]
$C_{\delta,\gamma} \lesssim c_{\delta,\gamma}^{-2}$,
uniform in $n \le \gamma N$ and $f \in C^1$
for fixed $\delta$, $\gamma$, $c$, $T$.
\end{remark}

\begin{corollary}[Closure of the bulk program]%
\label{cor:bulk-program-closed}
Theorem~\ref{thm:weak-limit} supplies hypothesis~(6.1) of
Lemma~6.3 of \cite{PaperIII} with
$\omega_n = C_{\delta,\gamma}/n$,
completing the logical chain:
\[
  \text{Thm.~\ref{thm:weak-limit}}
  \;\Rightarrow\; \texttt{lem:bulk-reduction}
  \;\Rightarrow\; \texttt{prob:comm-refined}
  \;\Rightarrow\; \text{Bulk tail bound}
\]
for all $n \le \gamma N$,
conditional on Assumption~2.4 of \cite{PaperIII}.
\end{corollary}

%%=================================================================
\section{Open Problems}\label{sec:open}
%%=================================================================

\begin{enumerate}[(P1)]
  \item \textbf{Extension to $f \in C^0$.}
  For $f \in C^0$, a density argument gives weak-$*$
  convergence without a rate.
  A rate for $f \in C^\alpha$ ($0 < \alpha < 1$)
  via interpolation between $C^0$ and $C^1$ would extend
  the result to H\"older-continuous test functions.

  \item \textbf{Off-diagonal regime.}
  An analogue for $|m-n| \ge \delta N$
  (cross terms $\psi_m\psi_n$) remains open
  and would complete the off-diagonal program of
  \cite{PaperIII}.

  \item \textbf{Explicit Bohr--Sommerfeld constant.}
  Making the constant in \eqref{eq:BS} fully explicit
  in $c$, $T$, $\gamma$ would allow the remainder in
  Theorem~\ref{thm:weak-limit} to be computed numerically.

  \item \textbf{Uniformity in $c$.}
  The regime $c \to \infty$ with $n/c$ fixed requires
  uniform WKB estimates in $c$ not directly available
  from \cite{OsipovRokhlinXiao}.

  \item \textbf{Near-critical regime $n \sim \gamma N$.}
  As $n \to \gamma N$, $I_\delta$ shrinks and the argument
  degenerates. An Airy-function scaling analysis is likely
  required near the spectral edge.
\end{enumerate}

%%=================================================================
\begin{thebibliography}{99}

\bibitem{PaperI}
Anonymous Author,
\emph{Coercivity of the Prolate Gram Form at Discrete
Sampling Points}, preprint, 2026.

\bibitem{PaperII}
Anonymous Author,
\emph{Scaling Limits of Prolate Gram Forms, Uniform
Coercivity, and a Trace Formula Toward Weil Positivity},
preprint, 2026.

\bibitem{PaperIII}
Anonymous Author,
\emph{PSWF Product Spectral Tail Estimates: Bulk and
Off-Diagonal Regimes, and an Obstruction at the
Spectral Edge}, preprint, 2026.

\bibitem{Bellman1943}
R.~Bellman,
\textit{The stability of solutions of linear differential
equations},
Duke Math.\ J.\ \textbf{10} (1943), 643--647.

\bibitem{Levitan}
B.~M.~Levitan and I.~S.~Sargsjan,
\textit{Sturm--Liouville and Dirac Operators},
Kluwer, Dordrecht, 1991.

\bibitem{OsipovRokhlinXiao}
A.~Osipov, V.~Rokhlin, and H.~Xiao,
\textit{Prolate Spheroidal Wave Functions of Order Zero},
Springer, New York, 2013.

\bibitem{Slepian1978}
D.~Slepian,
\textit{Prolate spheroidal wave functions, Fourier
analysis, and uncertainty~-- V},
Bell Syst.\ Tech.\ J.\ \textbf{57} (1978), 1371--1430.

\bibitem{Szego}
G.~Szeg\H{o},
\textit{Orthogonal Polynomials},
AMS Colloquium Publications, vol.~23, 1939.

\bibitem{Zygmund}
A.~Zygmund,
\textit{Trigonometric Series},
Cambridge University Press, 3rd~ed., 2002.

\end{thebibliography}

\end{document}
